%------------------------------------------------------------------------------
% Supplementary Material for: The Dimensional Landauer Bound
%------------------------------------------------------------------------------
\documentclass[aps,prx,twocolumn,superscriptaddress,floatfix]{revtex4-2}
\usepackage{amsmath,amssymb,amsfonts,graphicx,bm}
\usepackage{hyperref}

\begin{document}

\title{Supplementary Material: The Dimensional Landauer Bound}
\author{Ian Todd}
\affiliation{Sydney Medical School, University of Sydney, Sydney, NSW 2006, Australia}

\maketitle

%==============================================================================
\appendix
%==============================================================================

%------------------------------------------------------------------------------
\section{Formation Cost vs.\ Maintenance Cost}
\label{app:formation-maintenance}
%------------------------------------------------------------------------------

The Dimensional Landauer Bound (Eq.~3 in main text) establishes a minimum \emph{formation cost}: the thermodynamic work required to project a high-dimensional distribution onto a lower-dimensional manifold. This is analogous to Landauer's bound on bit erasure---a one-time energetic cost for a state transition.

However, biological and artificial systems typically operate as non-equilibrium steady states (NESS), continuously maintaining low-dimensional representations against thermal fluctuations. This incurs a \emph{maintenance cost}: ongoing power dissipation to sustain the projection.

The key result is that the geometric complexity $C_\Phi$ governs both costs:
\begin{itemize}
\item \textbf{Formation:} $W_{\mathrm{form}} \geq k_B T\, C_\Phi$ (Joules, one-time)
\item \textbf{Maintenance:} $P_{\mathrm{maint}} \propto k_B T\, C_\Phi / \tau_{\mathrm{relax}}$ (Watts, continuous)
\end{itemize}
where $\tau_{\mathrm{relax}}$ is the relaxation time of the confining potential.

The simulations in the main text measure maintenance power, which is the physically relevant quantity for living systems that must ``pay rent'' to sustain function. The scaling of maintenance power with geometric parameters validates the framework even though the simulations do not directly measure the formation bound.

%------------------------------------------------------------------------------
\section{Derivation of the Dimensional Landauer Bound}
\label{app:derivation}
%------------------------------------------------------------------------------

We derive the Dimensional Landauer Bound from stochastic thermodynamics, showing that the geometric contraction cost $C_\Phi$ arises necessarily from entropy production during dimensional projection.

\subsection{Setup: Projection as a physical process}

Consider a system with microstate $x \in \mathbb{R}^{D_\mathrm{in}}$ coupled to a heat bath at temperature $T$. The system evolves under overdamped Langevin dynamics:
\begin{equation}
dx_t = \mu F(x_t, \lambda_t) dt + \sqrt{2\mu k_B T}\, dW_t,
\end{equation}
where $\mu$ is mobility, $F$ is the total force (including control), $\lambda_t$ is a time-dependent protocol, and $W_t$ is a Wiener process.

A dimensional projection $\Phi: \mathbb{R}^{D_\mathrm{in}} \to \mathbb{R}^{D_\mathrm{out}}$ with $D_\mathrm{out} < D_\mathrm{in}$ is implemented by a control protocol that:
\begin{enumerate}
\item Confines the distribution to a neighborhood of the target manifold $\mathcal{M} = \Phi^{-1}(y)$
\item Maintains this confinement against thermal fluctuations
\end{enumerate}

\subsection{Entropy production during projection}

The total entropy production during a process from time $0$ to $\tau$ is~\cite{seifert2012}:
\begin{equation}
\Delta S_\mathrm{tot} = \Delta S_\mathrm{sys} + \Delta S_\mathrm{bath} \geq 0,
\end{equation}
where $\Delta S_\mathrm{bath} = Q/T$ and $Q$ is heat dissipated to the bath.

For the system entropy, we have:
\begin{equation}
\Delta S_\mathrm{sys} = S[p_\tau] - S[p_0] = -k_B \int p_\tau \ln p_\tau + k_B \int p_0 \ln p_0.
\end{equation}

\subsection{The projection bound}

The key insight is that dimensional projection changes the \emph{support} of the distribution. Let $p_0(x)$ be the initial distribution over $\mathbb{R}^{D_\mathrm{in}}$ and $p_\tau(y)$ be the final distribution over $\mathbb{R}^{D_\mathrm{out}}$ after projection.

The minimal entropy production for this transformation is bounded by the relative entropy between the initial distribution and any distribution that could have produced the final state via the projection $\Phi$:
\begin{equation}
\Delta S_\mathrm{tot} \geq k_B D_\mathrm{KL}(p_0(x) \| \Phi^\dagger p_\tau(y)),
\end{equation}
where $\Phi^\dagger p_\tau(y)$ is the lifted distribution---the maximum entropy distribution on $\mathbb{R}^{D_\mathrm{in}}$ consistent with the marginal $p_\tau(y)$ under $\Phi$.

This follows from the data processing inequality: the KL divergence cannot increase under the application of a stochastic map, and the projection $\Phi$ is precisely such a map.

\subsection{Work-entropy relation}

From the first law, the work done on the system equals:
\begin{equation}
W = \Delta F + Q,
\end{equation}
where $\Delta F$ is the change in free energy. For an isothermal process:
\begin{equation}
W = \Delta F + T \Delta S_\mathrm{bath} = \Delta F + T(\Delta S_\mathrm{tot} - \Delta S_\mathrm{sys}).
\end{equation}

For a projection that reduces information content by $\Delta I$ bits while incurring geometric contraction cost $C_\Phi$, the minimum work satisfies:
\begin{equation}
\boxed{W_\mathrm{min} \geq k_B T \ln 2 \cdot \Delta I + k_B T \cdot C_\Phi,}
\end{equation}
where:
\begin{itemize}
\item $k_B T \ln 2 \cdot \Delta I$ is the standard Landauer cost of erasing $\Delta I$ bits
\item $k_B T \cdot C_\Phi$ is the \emph{dimensional work}---the geometric cost of the projection itself
\end{itemize}

\subsection{Geometric interpretation of $C_\Phi$}

The contraction cost has a geometric interpretation. For a smooth projection $\Phi$ with Jacobian $J_\Phi$, the infinitesimal volume element transforms as:
\begin{equation}
dV_\mathrm{out} = \sqrt{\det(J_\Phi J_\Phi^\top)}\, dV_\mathrm{in}.
\end{equation}

The contraction cost is the expected log-volume ratio:
\begin{equation}
C_\Phi = \left\langle \ln \frac{dV_\mathrm{in}}{dV_\mathrm{out}} \right\rangle = -\frac{1}{2}\left\langle \ln \det(J_\Phi J_\Phi^\top) \right\rangle_{p_0(x)}.
\end{equation}

For a projection onto a curved manifold with principal curvatures $\kappa_i$, this includes contributions from both:
\begin{enumerate}
\item \textbf{Dimensional reduction}: $\frac{1}{2}(D_\mathrm{in} - D_\mathrm{out})\ln(2\pi e \sigma^2)$ for Gaussian fluctuations of width $\sigma$
\item \textbf{Curvature}: Additional terms proportional to $\kappa_i^2 \sigma^2$ from the non-Euclidean geometry of the target manifold
\end{enumerate}

%------------------------------------------------------------------------------
\section{Relation to Rate-Distortion Theory}
\label{app:rate-distortion}
%------------------------------------------------------------------------------

Shannon's rate-distortion theory~\cite{cover2006} establishes the minimum number of bits required to represent a source $X$ with expected distortion at most $D$:
\begin{equation}
R(D) = \min_{p(\hat{x}|x): \mathbb{E}[d(X,\hat{X})] \leq D} I(X; \hat{X}).
\end{equation}

The Dimensional Landauer Bound is \emph{complementary} to rate-distortion theory:

\begin{center}
\begin{tabular}{lcc}
\hline
& \textbf{Rate-Distortion} & \textbf{Dimensional Landauer} \\
\hline
Question & How many bits? & How many joules? \\
Units & bits & $k_B T$ (energy) \\
Constraint & Distortion $D$ & Fidelity $\varepsilon$ \\
Output & $R(D)$ & $W_\mathrm{min}$ \\
\hline
\end{tabular}
\end{center}

Rate-distortion theory tells you the minimum \emph{information} required for lossy compression. The Dimensional Landauer Bound tells you the minimum \emph{energy} required to physically implement that compression in a thermal environment.

The connection is:
\begin{equation}
W_\mathrm{min} \geq k_B T \ln 2 \cdot R(D) + k_B T \cdot C_\Phi(D),
\end{equation}
where $C_\Phi(D)$ is the geometric cost of the optimal encoder achieving rate $R(D)$.

Crucially, even when $R(D) = 0$ (lossless transmission of already-compressed data), dimensional work $C_\Phi > 0$ may still be required if the physical representation changes geometry.

%------------------------------------------------------------------------------
\section{Quantitative Validation of the Bound}
\label{app:validation}
%------------------------------------------------------------------------------

We validate the Dimensional Landauer Bound quantitatively in our simulations.

\subsection{Simulation 1: Curved manifold confinement}

For Brownian motion confined to the manifold $x_2 = A\sin(Bx_1)$, the theoretical contraction cost is:
\begin{equation}
C_\Phi^\mathrm{theory} = \frac{1}{2}\ln(1 + A^2 B^2 \langle\cos^2(Bx_1)\rangle) + \frac{k\sigma_\perp^2}{2k_B T},
\end{equation}
where $\sigma_\perp^2$ is the variance of fluctuations perpendicular to the manifold and $k$ is the confinement stiffness.

For our parameters ($A=1$, $B=2$, $k=10$, $D=0.1$, $T=1$):
\begin{itemize}
\item Theoretical: $C_\Phi^\mathrm{theory} \approx 0.69$
\item Measured: $\langle W_\mathrm{ctrl}/T \rangle / k_B T \approx 0.75 \pm 0.08$
\end{itemize}

The measured work exceeds the theoretical bound, as required by the second law, with the excess representing irreversible losses.

\subsection{Simulation 2: Kuramoto oscillators}

For $N$ coupled oscillators with order parameter $r$, the effective dimensionality is:
\begin{equation}
D_\mathrm{eff} = \frac{(\sum_i \lambda_i)^2}{\sum_i \lambda_i^2} \approx 1 + (N-1)(1-r^2),
\end{equation}
where $\lambda_i$ are eigenvalues of the phase covariance matrix.

The contraction cost for driving a 1D readout from $N$ oscillators scales as:
\begin{equation}
C_\Phi \sim \ln D_\mathrm{eff} \sim \ln N - r^2 \ln N.
\end{equation}

This predicts control work should decrease with $r^2$, which we observe (Fig.~2).

\subsection{Simulation 3: Autoencoder bottleneck}

For data with intrinsic dimension $d^*$ compressed to bottleneck dimension $d_b$, the contraction cost diverges as:
\begin{equation}
C_\Phi \sim \begin{cases}
O(1) & d_b \geq d^* \\
O((d^* - d_b)\ln(1/\varepsilon)) & d_b < d^*
\end{cases}
\end{equation}
where $\varepsilon$ is reconstruction error.

This explains the sharp increase in training effort (SGD work proxy) below $d_b = d^* = 2$ (Fig.~3).

%------------------------------------------------------------------------------
\section{Experimental Proposals}
\label{app:experiments}
%------------------------------------------------------------------------------

We propose three experimental tests of the Dimensional Landauer Bound.

\subsection{Colloidal particles in optical traps}

\textbf{Setup}: A colloidal particle in a 2D optical trap, with feedback control confining it to a 1D curved path.

\textbf{Protocol}:
\begin{enumerate}
\item Initialize particle in 2D Gaussian equilibrium
\item Apply feedback control to confine to curve $y = f(x)$ with varying curvature
\item Measure heat dissipation via position tracking and Jarzynski equality
\item Compare to theoretical $W_\mathrm{dim} = k_B T C_\Phi$
\end{enumerate}

\textbf{Prediction}: Heat dissipation should scale with integrated squared curvature $\int \kappa^2(s) ds$.

\textbf{Feasibility}: Standard optical tweezer setup with video microscopy. Similar to Bérut et al.~\cite{berut2012} Landauer verification.

\subsection{Electronic feedback circuits}

\textbf{Setup}: An RC circuit with $N$ noisy nodes, feedback-controlled to produce a 1D output voltage.

\textbf{Protocol}:
\begin{enumerate}
\item Measure thermal noise spectrum across all nodes
\item Apply linear feedback to project onto target 1D subspace
\item Measure power dissipation in feedback amplifiers
\item Vary projection geometry and compare to theory
\end{enumerate}

\textbf{Prediction}: Dissipation $\propto k_B T \ln(N)$ for random projection, reduced for coherent (eigenmode-aligned) projection.

\subsection{Neural population recordings}

\textbf{Setup}: Multi-electrode array recordings from motor cortex during reaching tasks.

\textbf{Protocol}:
\begin{enumerate}
\item Estimate neural manifold dimensionality $D_\mathrm{eff}$ via participation ratio
\item Quantify metabolic cost via BOLD or direct oxygen measurements
\item Correlate $D_\mathrm{eff}$ with metabolic rate across conditions
\end{enumerate}

\textbf{Prediction}: Tasks requiring low-dimensional readout from high-dimensional population should show elevated metabolic cost scaling as $\sim \ln(D_\mathrm{eff}/D_\mathrm{out})$.

%------------------------------------------------------------------------------
\section{Connection to Existing Bounds}
\label{app:connections}
%------------------------------------------------------------------------------

The Dimensional Landauer Bound connects to several established results:

\textbf{Sagawa-Ueda bound}~\cite{sagawa2010}: For feedback-controlled systems, the extractable work is bounded by mutual information. Our bound is complementary: the work \emph{required} for dimensional projection is bounded by the contraction cost.

\textbf{Jarzynski equality}~\cite{jarzynski1997}: Our bound is consistent with $\langle e^{-W/k_B T}\rangle = e^{-\Delta F/k_B T}$. The dimensional work appears in the non-equilibrium free energy difference.

\textbf{Crooks fluctuation theorem}~\cite{crooks1999}: The bound can be tightened using Crooks' relation to account for fluctuations in the projection process.

\textbf{Information geometry}~\cite{amari2016}: The contraction cost $C_\Phi$ is related to the Fisher information metric on the statistical manifold. Projections that preserve Fisher information incur lower dimensional work.

\bibliography{references}

\end{document}
